\section{What a Harness Is}
\label{sec:harness}

\paragraph{The model.} We take a language model to be a function from a token sequence to a distribution over the next token, applied autoregressively until a stop token is produced:
\[
  \textit{TokensOut} = \textit{Model}(\textit{TokensIn}).
\]
Chat templates, tool-call formats and reasoning tokens are conventions for laying out \textit{TokensIn} and reading \textit{TokensOut}; they are learned by the model in training, but they are written and parsed by something else.

\paragraph{The harness.} We define the \emph{harness} as everything in an agent that is not the model: the software that turns the world into tokens for the model, and the model's tokens into effects on the world. An agent is a model inside a harness:
\[
  \textit{Agent} = \textit{Harness}(\textit{Model}).
\]
The definition draws the boundary at the model's weights and its sampling. Everything the model is shown, everything it is able to do, and everything that happens between one call and the next is on the harness side. This includes things often attributed to ``the model'', such as memory across sessions (the model has none; the harness stores and retrieves it), tool use (the model emits a request; the harness carries it out), and knowledge of the current date or working directory (the model knows only what the harness writes into its context).

\paragraph{The loop.} Figure~\ref{fig:loop} shows the shape every harness shares. Something comes in: a request from a person, an event, or the result of an earlier action. The harness assembles a context around it and sends that context to the model. The model runs its own inner loop, one token at a time, until it stops. The harness handles the response: showing words to a person, or carrying out the actions it requests. If an action produced a result, the result comes back in, and the loop runs again; if not, control returns to the person, or the agent stops.

\begin{figure}[t]
\centering
\resizebox{\linewidth}{!}{%
\begin{tikzpicture}[
  node distance=8mm and 8mm,
  box/.style={draw, rounded corners=2pt, minimum width=21mm, minimum height=9mm, align=center, font=\small},
  model/.style={box, fill=black!85, text=white, minimum width=27mm},
  arr/.style={-{Stealth[length=2mm]}, thick}]
  \node[box] (in) {\prim{Input}\\[-1pt]\scriptsize perceive};
  \node[box, right=of in] (ctx) {\prim{Context}\\[-1pt]\scriptsize working memory};
  \node[model, right=17mm of ctx] (m) {Model\\[-1pt]\scriptsize tokens $\to$ next token};
  \node[box, right=of m] (out) {\prim{Output}\\[-1pt]\scriptsize act};
  \node[font=\small, right=10mm of out] (p) {person};
  \draw[arr] (in) -- (ctx);
  \draw[arr] (ctx) -- node[above, font=\scriptsize]{\prim{Model}} node[below, font=\scriptsize]{\prim{interface}} (m);
  \draw[arr] (m) -- (out);
  \draw[arr] (out) -- (p);
  \draw[arr] (out.south) -- ++(0,-8mm) -| node[pos=0.25, below, font=\scriptsize]{a tool's result comes back in: go again} (in.south);
  \draw[arr, dashed] (p.north) -- ++(0,7mm) -| node[pos=0.25, above, font=\scriptsize]{a request} (in.north);
  \node[draw, dotted, rounded corners=4pt, inner sep=4mm, fit=(in)(ctx)(m)(out), label={[font=\scriptsize]north west:\prim{Control flow} holds the loop}] {};
\end{tikzpicture}}
\caption{The agent loop. The harness's five primitives surround the model; \prim{Control flow} decides whether the loop runs again, hands back, or stops.}
\label{fig:loop}
\end{figure}

\paragraph{Not every harness is an agent's.} A program that makes one model call and prints the reply has a harness, but a trivial one: its input is fixed, its context is one message, its output is printed, and its control flow is a straight line. The five primitives describe this case as well; what turns a harness into an \emph{agent's} harness is a \prim{Control flow} that feeds the results of actions back to the model (\S\ref{sec:primitives}).
